\documentclass[11pt,letterpaper]{article}

\usepackage[utf8]{inputenc}
\usepackage[spanish]{babel}
\usepackage[margin=2.5cm]{geometry}
\usepackage{amsmath,amssymb}
\usepackage{enumitem}
\usepackage{titlesec}
\usepackage{array}
\usepackage{hyperref}
\usepackage{xcolor}

\definecolor{unicaucaverde}{RGB}{0,90,60}
\hypersetup{colorlinks=true,linkcolor=unicaucaverde,urlcolor=unicaucaverde}

\titleformat{\section}{\normalfont\Large\bfseries\color{unicaucaverde}}{\thesection}{1em}{}
\titleformat{\subsection}{\normalfont\large\bfseries}{\thesubsection}{1em}{}

\setlength{\parindent}{0pt}
\setlength{\parskip}{6pt}

\title{\textcolor{unicaucaverde}{\Large Guía de Aprendizaje --- Semana 4}\\[4pt]
  \large Ámbito avanzado y recursión}
\author{Programación --- Universidad del Cauca}
\date{2026}

\begin{document}
\maketitle

\begin{center}
\begin{tabular}{|>{\bfseries}p{4cm}|p{10cm}|}
\hline
Semana & 4 de 16 \\ \hline
Unidad & Fundamentos de programación en C \\ \hline
Subtemas & Variables estáticas locales y tiempo de vida (ámbito avanzado), funciones recursivas
  (caso base, caso recursivo, pila de llamadas) \\ \hline
Horas de clase & 4 \\ \hline
Resultado de aprendizaje & RA3 (Módulo 3, parte 2 de 2; continúa de la Semana 3, que cubrió
  funciones y ámbito básico) \\ \hline
\end{tabular}
\end{center}

\section*{Relevancia para ingeniería}

Muchos algoritmos de ingeniería tienen una estructura naturalmente recursiva: un cálculo numérico
iterativo que refina una aproximación paso a paso, un recorrido de árbol de decisiones, o una
búsqueda que divide un problema en subproblemas idénticos más pequeños. Reconocer cuándo un
problema se resuelve de forma más clara con recursión --- y entender su costo en memoria a través
de la pila de llamadas --- es una herramienta más en el diseño de soluciones, junto con las
variables estáticas locales, que permiten que una función recuerde estado entre llamadas sin
recurrir a una variable global.

\section{Objetivo de la semana}

Distinguir tiempo de vida de ámbito usando variables estáticas locales, y diseñar funciones
recursivas correctas (con caso base y caso recursivo) para resolver problemas naturalmente
recursivos, comparando su comportamiento con una solución iterativa equivalente.

\section{Resultados de aprendizaje de la semana}

\begin{itemize}[leftmargin=1.2cm]
  \item[\textbf{RA3.3.}] Explicar la diferencia entre ámbito y tiempo de vida de una variable, y
        usar variables estáticas locales (\texttt{static}) para que una función mantenga estado
        entre llamadas sucesivas.
  \item[\textbf{RA3.4.}] Diseñar y escribir funciones recursivas simples, identificando su caso
        base y su caso recursivo, explicar su ejecución mediante la pila de llamadas, y compararlas
        con una solución iterativa equivalente.
\end{itemize}

\section{Contenidos}

\begin{itemize}[leftmargin=1.2cm]
  \item \textbf{4.1.} Variables estáticas locales (\texttt{static}): tiempo de vida frente a
        ámbito, bloques anidados y sombreado.
  \item \textbf{4.2.} Recursión: caso base, caso recursivo, pila de llamadas, recursión frente a
        iteración.
\end{itemize}

\section{Plan de la sesión (4 horas)}

\begin{enumerate}[leftmargin=1.2cm]
  \item[\textbf{Hora 1.}] Control formativo corto (10 min) de la Semana 3 (funciones, paso por
        valor, ámbito local/global); qué significa \texttt{static} en una variable local, tiempo de
        vida frente a ámbito (teoría + ejemplos).
  \item[\textbf{Hora 2.}] Taller: funciones con contador interno usando \texttt{static}
        (\texttt{contadorLlamadas}, \texttt{siguienteId}); trazar a mano el valor de una variable
        estática tras varias llamadas.
  \item[\textbf{Hora 3.}] Recursión: qué es, caso base y caso recursivo, la pila de llamadas
        (teoría + ejemplo \texttt{factorial}).
  \item[\textbf{Hora 4.}] Taller: escribir funciones recursivas (\texttt{potencia},
        \texttt{sumaHasta}, \texttt{fibonacci}) y comparar cada una con su versión iterativa
        equivalente.
\end{enumerate}

\section{Actividades}

\begin{itemize}[leftmargin=1.2cm]
  \item \textbf{En clase}: talleres de las horas 2 y 4 (ver arriba).
  \item \textbf{Trabajo autónomo}: ejercicios propuestos de \texttt{notas.tex} §5.
  \item \textbf{Software}: un compilador de C (\texttt{gcc} local o compilador en línea) para
        escribir, compilar y ejecutar cada ejemplo y ejercicio.
\end{itemize}

\section{Material de apoyo}

\begin{itemize}[leftmargin=1.2cm]
  \item \texttt{notas.tex} --- desarrollo teórico completo de 4.1 y 4.2, con ejemplos de código y
        ejercicios resueltos.
  \item \texttt{diapositivas.tex} --- síntesis visual de la sesión.
\end{itemize}

\section{Evaluación de la semana}

Taller calificado al cierre de la Hora 4: escribir al menos dos funciones recursivas (una de ellas
comparada explícitamente con su versión iterativa) y una función con variable estática local que
mantenga estado correcto entre llamadas, para verificar RA3.3 y RA3.4 antes de avanzar a la Semana
5 (arreglos unidimensionales).

\section{Bibliografía de la semana}

\begin{enumerate}[leftmargin=1.2cm]
  \item Deitel, P., Deitel, H. \emph{C: How to Program}. 8a ed., Pearson, cap. 5 (funciones
        recursivas y clase de almacenamiento \texttt{static}).
  \item Kernighan, B. W., Ritchie, D. M. \emph{The C Programming Language}. 2a ed., Prentice Hall,
        cap. 4.
  \item Cairo, L. J. \emph{Fundamentos de programación}. 4a ed., McGraw-Hill.
\end{enumerate}

\end{document}
